\clearpage
\setcounter{page}{1}
\maketitlesupplementary

\section{Dataset Statistics}
\label{sec:suppl-dataset}

\begin{table}[h]
  \centering
  \caption{Dataset statistics. The LoRA split uses a scene-level partition with no
  scene overlap between train, val, and test.}
  \label{tab:dataset_stats}
  \begin{tabular}{lrrr}
    \toprule
    & \textbf{Train} & \textbf{Val} & \textbf{Test} \\
    \midrule
    Scenes      & 20    & 7   & 20   \\
    Questions   & 2535  & 983 & 2706 \\
    \midrule
    \multicolumn{4}{l}{\small Total: 47 scenes, 6224 questions} \\
    \bottomrule
  \end{tabular}
\end{table}

\section{LoRA Fine-Tuning Details}
\label{sec:suppl-lora}

\begin{table}[h]
  \centering
  \caption{LoRA fine-tuning hyperparameters for Qwen3-4B.}
  \label{tab:lora_hparams}
  \begin{tabular}{ll}
    \toprule
    \textbf{Hyperparameter} & \textbf{Value} \\
    \midrule
    LoRA rank               & 16 \\
    LoRA alpha              & 32 \\
    LoRA dropout            & 0.05 \\
    Target modules          & q, k, v, o projections \\
    Epochs                  & 3 \\
    Learning rate           & $2\times10^{-4}$ \\
    Batch size              & 1 (grad.\ accum.\ 8; effective 8) \\
    Max sequence length     & 1024 tokens \\
    Seed                    & 42 \\
    \midrule
    Train scenes / questions  & 20 / 2535 \\
    Val scenes / questions    & 7 / 983 \\
    Test scenes / questions   & 20 / 2706 \\
    \bottomrule
  \end{tabular}
\end{table}

The train/val/test split is scene-level with no overlap, preventing the model from
memorizing specific visual scenes. We compare against zero-shot Qwen3-4B on the identical
test split, giving an apples-to-apples BASE vs.\ +LoRA comparison.
